% ============================================================================
\clearpage
\section*{v4.10.0-r06 Source-Only: Semiclassical Einstein Equations from Relative Entropy (Dorau--Much) and the Readouts of the Companion}
\addcontentsline{toc}{section}{v4.10.0-r06 Source-Only: Semiclassical Einstein Equations from Relative Entropy (Dorau--Much) and the Readouts of the Companion}
\label{sec:v4100-r06-relative-entropy}
% ============================================================================

\paragraph{Purpose.}
This layer takes up a recent result from the literature, pointed out by the author: P.~Dorau and A.~Much (Leipzig), \emph{From Quantum
Relative Entropy to the Semiclassical Einstein Equations}, Phys.\ Rev.\ Lett.\ \textbf{136}, 091602 (2026), DOI
\href{https://doi.org/10.1103/lmq8-nsty}{10.1103/lmq8-nsty}, arXiv:2510.24491.  It records what the paper shows and assumes, and where it meets the
readouts of the companion: the entropy factor of \S\ref{sec:bh} (v4.8.0-r03 F11a), the free \(\Lambda\) (v4.10.0-r03, r05), the sign of gravity
(v4.10.0-r01) and the frames.  The layer is an exploration of the literature; it contains no pre-registered test and no new prediction.
The note is \texttt{notes/\allowbreak research/\allowbreak v4\_10\_0\_r06/}.  Earlier layers are not rewritten.

\noindent\textbf{Source-only status.}
This v4.10.0-r06 layer is source-only.  It does not build a PDF, does not assign a new DOI,
and does not authorize a public release.  The v4.10.0 Version DOI
\href{https://doi.org/10.5281/zenodo.23991357}{10.5281/zenodo.23991357} and the Concept DOI
\href{https://doi.org/10.5281/zenodo.19210728}{10.5281/zenodo.19210728} are unchanged by it.

\subsection*{v4.10.0-r06 results}
\label{subsec:v4100-r06-results}

\paragraph{Theoremlet v4.10.0-r06.1 (the result of Dorau and Much, as stated by them).}
Literature, read in arXiv version v3; the formulas were not recomputed here.
\begin{enumerate}[leftmargin=2em]
\item[(a)] \emph{Setting.} Near a point the metric is approximated by \(\eta_{ab}\) (leading order).  The boost field \(\xi\) is null on two null
hypersurfaces meeting in a spacelike surface \(\mathcal S\): a local Rindler horizon, as in Jacobson's thermodynamic derivation (1995).  A real
Klein--Gordon field of mass \(m\) is quantized as a CCR algebra.  The reference state \(\omega_0\) is quasifree, boost invariant and Hadamard (in the
Rindler case the Minkowski vacuum); only coherent excitations \(\omega_\phi\) are treated.
\item[(b)] \emph{Relative entropy equals energy flux.} The modular group of the wedge acts as the boost, \(\Delta^{it}=\) boost by \(2\pi t\)
(Bisognano--Wichmann, Summers--Verch), and \(\omega_0\) is KMS at \(\beta=2\pi/\kappa\).  The Araki relative entropy then becomes
\[
S_{\rm rel}(\omega_0\Vert\omega_\phi)=-2\pi\int U\,\langle{:}T_{ab}{:}\rangle_{\omega_\phi}\xi^a\xi^b\,dU\,d\mathrm{vol}_{\mathcal S}.
\]
\item[(c)] \emph{Geometry.} With vanishing expansion, shear and twist at the horizon, the Raychaudhuri equation gives
\(\delta A=-\int U\,R_{ab}\xi^a\xi^b\,dU\,d\mathrm{vol}_{\mathcal S}\).
\item[(d)] \emph{Assumed, not derived:} \(S_{\rm rel}=\delta A/4\) (Bekenstein--Hawking, units \(G=\hbar=1\)).  This fixes the coupling \(\alpha=8\pi\).
\item[(e)] \emph{Result.} Holding for all null directions, \(\alpha\langle T_{ab}\rangle=R_{ab}+Ng_{ab}\); conservation and the contracted Bianchi identity
give \(N=-R/2+\Lambda\), hence \(G_{ab}+\Lambda g_{ab}=8\pi\langle{:}T_{ab}{:}\rangle\) with \(\Lambda\) an arbitrary integration constant.
\end{enumerate}
The authors call this ``arguments indicating'' the result: leading order, coherent states, scalar field.  The press reading ``gravity from vacuum
energy'' is not what is shown: the vacuum enters as the reference state of a relative entropy, not as a source of energy.

\paragraph{Theoremlet v4.10.0-r06.2 (the entropy factor becomes the gravitational coupling).}
Analytic, not blind (follows from r06.1 by one line).  Let a horizon counting give \(S=c\,A/4\) with \(c>0\).  Replacing (d) by
\(S_{\rm rel}=c\,\delta A/4\) and keeping the rest of r06.1 gives
\[
\alpha=\frac{8\pi}{c},\qquad G_{\rm eff}=\frac{G}{c}.
\]
For the two countings of \S\ref{sec:bh} (F11a in \S\ref{subsec:v480-r03-frame-parameter-provenance}) one has \(c=\gamma_*/\gamma\), where \(\gamma_*\) is the value
at which the counting reproduces \(S_{\rm BH}\): \(\gamma_{\rm string}=\ln2/(\pi\sqrt3)\) for the counting with \(j=\tfrac12\) only, \(\gamma_{\rm BI}=0.237533\) for
the counting of Domagala--Lewandowski and Meissner.  At one common value of \(\gamma\) the two countings differ by the factor
\(\gamma_{\rm BI}/\gamma_{\rm string}=\gamma_L=1.864700\) (F11a).  Under the chain of r06.1 this factor would appear as a ratio of gravitational couplings:
\[
\frac{G_{\rm eff}^{(j=1/2)}}{G_{\rm eff}^{(\rm DL/M)}}=\frac{c_{\rm DL/M}}{c_{j=1/2}}=\frac{\gamma_{\rm BI}}{\gamma_{\rm string}}=\gamma_L .
\]
\emph{Consequence.} If the T frame counts with \(j=\tfrac12\) and the M frame with all spins (r63, H47) at one common \(\gamma\), the two frames see different
Newton constants.  A single measured \(G\) is consistent only if each frame uses the \(\gamma\) at which its own counting reproduces \(S_{\rm BH}\); then
\(c=1\) in both, \(G\) is frame independent, and the frame difference sits in \(\gamma\).  This is the reading of r64 (the frame difference in \(\gamma\), not in
\(\hbar\) or \(G\)); Dorau--Much make it a statement about the coupling, not about the entropy alone.

\paragraph{Theoremlet v4.10.0-r06.3 (\(\Lambda\) as integration constant).}
Literature and comparison.  In r06.1(e) \(\Lambda\) is not fixed by the derivation; it enters through \(\nabla^a(\Lambda g_{ab})=0\).  This is the same situation
as in v4.10.0-r03 and r05: under G2 the action gives no \(\Lambda\) after the bang, \(\Lambda=a\,V_{\min}\) is free, and \(\Lambda\) is a readout
(Theorem~\ref{thm:dark-energy}).  The agreement supports the classification ``\(\Lambda\) is not part of the dynamics''.  It is not a success of the companion:
the standard derivation does not fix \(\Lambda\) either.  The readout \(\Omega_M=\beta_{\rm sym}^2\) (H96) is neither supported nor contradicted.

\paragraph{Remark v4.10.0-r06.4 (the sign of gravity from positivity; afterwards).}
In r06.1 the sign of \(\alpha\) is not free: \(S_{\rm rel}\ge0\) and \(c>0\) give \(\alpha>0\), i.e.\ attractive gravity.  In v4.10.0-r01 the sign of W1 was
corrected by an author decision after H90 had shown repulsion.  The corrected sign coincides with the one fixed by positivity.  This is noted afterwards
and is not a derivation of W1: Dorau--Much use a scalar field and Riemannian geometry, W1 is Einstein--Cartan with spinors and torsion
\(+(3/16)\kappa j^2\).  A transfer would need the relative entropy for Dirac fields, which the paper does not treat.

\paragraph{Remark v4.10.0-r06.5 (frames, Wick rotation and the modular group; open).}
The derivation uses that the modular group of the reference state is a real one-parameter group of boosts with KMS period \(2\pi\).  In the Euclidean
picture that period is the full rotation angle \(2\pi\) (the Unruh temperature via Wick rotation).  The companion has this bridge itself
(Theorem~\ref{thm:artqft}, \(t\mapsto i\tau_E\) on the biquaternion).  Under the tachyon map \(z\mapsto iz\) of Convention~F the rapidity becomes
imaginary and the boost a rotation.  Whether \(z\mapsto iz\) coincides with the Wick rotation of Theorem~\ref{thm:artqft} is open.  If it does, frame T is
the Euclidean side of the same KMS structure, not its loss.  The precise question is operator algebraic: is there a state on the wedge algebra whose
modular group is the tachyon image of the boost?  No answer is given here.

\paragraph{Proof and measurement sources.}
Exploration note \texttt{Erkundung\_\allowbreak Dorau\_\allowbreak Much\_\allowbreak relative\_\allowbreak Entropie.md} in
\texttt{notes/\allowbreak research/\allowbreak v4\_10\_0\_r06/}; the paper (arXiv:2510.24491v3).  r06.2 is one line of algebra on r06.1(d)--(e).  \(\square\)

\paragraph{Consequence.}
Dorau and Much give a rigorous quantum field theoretic form of Jacobson's argument.  They assume the area law and leave \(\Lambda\) free.  For the
companion this has three uses: the entropy factor of \S\ref{sec:bh} is the gravitational coupling (r06.2), so the frame reading of the horizon
countings must keep \(G\) frame independent; \(\Lambda\) as a readout is consistent with the standard derivation (r06.3); and the frames get a precise
operator algebraic question (r06.5).  Neither \(G\) nor \(\Lambda\) follows from it.

\subsection*{v4.10.0-r06 status in the three frames}
\label{subsec:v4100-r06-frames}

\begin{itemize}[leftmargin=2em]
\item M: the derivation of r06.1 holds as stated (real boosts, KMS at \(2\pi/\kappa\)).
\item T: open (r06.5); under \(z\mapsto iz\) the boost becomes a rotation, and whether the KMS structure survives as its Euclidean side is not decided.
\item P: degenerate fixed point; no separate statement.
\end{itemize}
r06.2 is frame dependent by construction (one counting per frame); r06.3 holds in every frame that has the derivation.  No frame-free statement is made.

\subsection*{v4.10.0-r06 register entries}
\label{subsec:v4100-r06-register}

\begin{center}
\small
\begin{tabular}{@{}>{\raggedright\arraybackslash}p{0.34\textwidth}>{\raggedright\arraybackslash}p{0.58\textwidth}@{}}
\toprule
statement & status after v4.10.0-r06\\
\midrule
Einstein equations from relative entropy & literature (Dorau--Much 2026); area law assumed, \(\Lambda\) free\\
entropy factor of \S\ref{sec:bh} & equals the gravitational coupling (r06.2); \(G\) frame independent only with matched \(\gamma\) per frame\\
\(\Lambda\) & integration constant in the standard derivation; readout in the companion (consistent)\\
sign of W1 (v4.10.0-r01) & coincides with the positivity sign (afterwards, scalar case only)\\
frame T and the modular group & open question (r06.5)\\
\bottomrule
\end{tabular}
\end{center}

\subsection*{v4.10.0-r06 ledger and guard update}
\label{subsec:v4100-r06-ledger-guard}

\[
R_{\rm chk}^{v4100\text{-}r06}=(1,1,0)
\]
for (literature recorded; entropy factor tied to the coupling; frame T under the modular group open).

The layer spends no guard:
\[
R_{\rm glob}^{(v4100\text{-}r06,+)}=(0,1,1,1,1),
\qquad
\Gamma_{\rm guard}^{v4100\text{-}r06}=(1,1,1,1,1),
\qquad
\beta_4=1/256\ \text{unspent}.
\]
It does not imply exact finite \(C5\) zero closure, exact all-readout terminality,
\(C6/T2\), a global photon-clock/RH theorem, or global all-readout closure.
